\scanpage{483}
\mspage{245}
34. Description heuristique profinie de la catégorie
des courbes alg. \unsure{définies} sur \deleted{\unsure{$\mathbb Q$}} des sous-
extensions finies $K$ de $\overline{\mathbb Q}_0/\mathbb Q$ (i.e. de $\mathbb C/\mathbb Q$).

On se borne aux courbes \inserted{géom.} connexes \inserted{(par commodité)}, anabéliennes
(par nécessité provisoire \unsure{!}, cf plus bas sur \unsure{la façon} de
se débarrasser de cette restriction). La donnée d'
un rev. universel de \deleted{un} $U_{\overline{\mathbb Q}_0}$ définit une
structure d'extension
\[
(1)\qquad 1\longrightarrow\pi\longrightarrow\mathcal E\longrightarrow\Gamma^{+}\longrightarrow1\qquad\text{ou}
\]
(2) ou encore on a un hom $\mathcal E\longrightarrow\boldsymbol{\Gamma}$
(image $\Gamma^{+}$ ouverte dans $\boldsymbol{\Gamma}$,
noyau appelé $\pi$),
où $\Gamma^{+}\subset\boldsymbol{\Gamma}_{\mathbb Q}\simeq\boldsymbol{\Gamma}_{0,3}$ est le \unsure{ss-gr.|sous-gr.} ouvert
correspondant à $K$. Se donner une telle
extension (moyennant $\unsure{\mathrm{Centre}}(\pi)=1$) revient à
se donner une action extérieure de $\Gamma'$ sur $\pi$.
Une première question : faut-il mettre la
structure à lacets de $\pi$ dans les \deleted{\illegible{1}} données \unsure{$\mathcal E$} de
l'objet (1) (ou \unsure{(2)|(1')}) censé décrire $U\unsure{\mathcal E}/K$ ?

\uline{Conjecture} Ce n'est pas la peine -- la structure à lacets
de $\pi$ est la seule structure à lacets \unsure{invariante}
par l'action extérieure de $\Gamma^{+}$ (ou de $\Gamma^{\natural}$). Les
sous-groupes à lacets $L_i$ sont les sous-groupes
maximaux\inserted{dans $\pi$}, tels que $\mathrm{Norm}_{\mathcal E}(L_i)\longrightarrow\Gamma^{+}$ ait
comme image un ss-groupe ouvert de $\unsure{\Gamma^{+}|\Gamma}$.
\margin{\unsure{et} tout sous-groupe $\ne1$ de $\pi$ qui a cette propriété est contenu dans un unique $L_i\ldots$}
Je
présume aussi que tout hom \deleted{\unsure{d'}}\inserted{entre} extensions
$\mathcal E$ de $\Gamma^{+}$ par un $\pi$, $\mathcal E'$ de $\Gamma$ par un $\pi'$
($\mathcal E,\mathcal E'$ provenant de courbes algébriques $U,U'$) et tel
que l'image de $\mathcal E$ dans $\mathcal E'$ soit ouverte, respecte
\margin{482}
